% UBT-AI-PROVENANCE-BEGIN
% schema: ubt-ai-provenance/v1
% tier: B_machine_verified
% ai_assistance: disclosed
% human_review: machine-verification
% editorial_responsibility: Ing. David Jaroš
% policy: ../../AI_PROVENANCE.md
% notice: Machine-verified against named sources or verifiers; individual attestation is not claimed.
% UBT-AI-PROVENANCE-END

\documentclass[11pt,a4paper]{article}
\usepackage[margin=2.5cm]{geometry}
\usepackage{amsmath,amssymb,amsthm}
\usepackage[most]{tcolorbox}
\usepackage[hidelinks]{hyperref}
\usepackage{ubtprovenance}
\UBTTier{B}
\newtheorem{theorem}{Věta}
\newtheorem{corollary}[theorem]{Důsledek}
\title{Audit nepodmíněné reprodukce GR v kanonické UBT}
\author{Ing. David Jaro\v{s}}
\date{23. srpna 2026}
\begin{document}\maketitle
\UBTProvenanceNotice

% BILINGUAL-UNIT: unconditional.question
\section{Rozhodovací otázka}
Tento audit zkoumá, zda lze status lokální čtyřrozměrné infračervené
reprodukce GR změnit z \texttt{CLOSED\_CONDITIONALLY} na \texttt{CLOSED} při
použití současných kanonických axiomů a dokázaných výsledků repozitáře. Testuje
mikroskopický jedno-$\Theta$ selector, koeficient Einsteinova--Hilbertova členu
a míru, kompozitní Hessián, fyzikální stabilitu $\psi$-sektoru a logickou
relevanci globálního slepování pro lokální tvrzení.

% BILINGUAL-UNIT: unconditional.selector
\section{Přesná překážka selektoru}
Na Lorentzově reálné větvi nechť
\begin{equation}
 E_\mu=\mathcal N_0^{-1/2}D_\mu\Theta,
 \qquad \langle E_\mu,E_\nu\rangle_\sharp=g_{\mu\nu}.
\end{equation}
Potom nejpřímější kovariantní kvadratický kandidát z prvního jetu splňuje
\begin{equation}\label{eq:collapse}
 g^{\mu\nu}\langle D_\mu\Theta,D_\nu\Theta\rangle_\sharp
 =\mathcal N_0 g^{\mu\nu}g_{\mu\nu}=4\mathcal N_0.
\end{equation}
Jeho akce je tedy kosmologickým objemovým funkcionálem. Není nezávislým
kinetickým selektorem a nemůže odvodit Einsteinův--Hilbertův člen.

\begin{theorem}[Nepodmíněné uzavření nelze odvodit ze současných vstupů]
Nechť $\mathcal K$ je současná množina kinematických vstupů: definice
kovariantního tetrádu, centrální metrická identita, lokální Lorentzova a
difeomorfní symetrie, rekonstrukce konexe a lokální split-jet right inverse.
Pro každé reálné $c$ má funkcionál
\begin{equation}\label{eq:family}
 S_c[\Theta]=S_0[\Theta]+c\int d^4x\sqrt{-g[\Theta]}
 R[g[\Theta]]
\end{equation}
stejné vstupy $\mathcal K$. Různá $c$ dávají různé Newtonovy koeficienty,
včetně $c=0$. Proto $\mathcal K$ nevybírá Einsteinův--Hilbertův člen ani
neurčuje konečný kladný koeficient. Rovnice~\eqref{eq:collapse} tuto volnost
neodstraňuje. Předpoklady podmíněné věty o reprodukci tedy nejsou důsledky
aktuálně finalizované mikroskopické jedno-$\Theta$ teorie.
\end{theorem}
\begin{proof}
Každý člen rodiny v rovnici~\eqref{eq:family} má stejnou mapu pole na tetrád a
stejné kinematické identity. Přidaný člen je samostatně difeomorfně a lokálně
Lorentzovsky invariantní pro každé $c$, zatímco jeho metrická variace a
Newtonův koeficient závisejí na $c$. Společné vstupy proto nemohou implikovat
jednu hodnotu $c$ ani vyloučit $c=0$. Rovnice~\eqref{eq:collapse} ukazuje, že
uvedený kvadratický kandidát z prvního jetu přidává pouze objemový člen, a
nedodává tedy chybějící selekční premisu.
\end{proof}

% BILINGUAL-UNIT: unconditional.hessian
\section{Hessián, míra a $\psi$-sektor}
Výpočet indukované gravitace je přesný pouze po předpokladu kalibračně
fixovaného fyzikálního Hessiánu Laplaceova typu, počtu módů $N_B$, křivostní
vazby $\xi$, identifikace cutoffu, regulátoru a omezené fluktuační míry.
Hlavní symbol na fixovaném pozadí není úplný kompozitní Hessián: variace
$\Theta$ mění také $E[\Theta]$, $g[\Theta]$ a rekonstruovanou konexi. Žádná ze
současných vět tato chybějící data neodvozuje.

Stejně tak je čtyřrozměrná truncace na nulový mód fyzikálně vybrána pouze
tehdy, když ji finalizovaný mikroskopický vývoj zachovává a nežádoucí
$\psi$-módy jsou stabilní nebo se oddělují. Existující věta o propagaci
symetrie je podmíněna ekvivariancí a jednoznačným vývojem; tyto hypotézy nebyly
dokázány pro finalizovanou mikroskopickou akci. Stabilita $\psi$ tedy není
nutná k variaci již předpokládané čtyřrozměrné efektivní akce, je však nutná k
nepodmíněnému odvození této fyzikální větve z UBT.

% BILINGUAL-UNIT: unconditional.global
\section{Lokální versus globální rozsah}
Globální pokračování přes horizonty a nulové patche ani topologické
pokračování není premisou uvedené \emph{lokální} věty o reprodukci GR. Samo
proto nemůže blokovat lokální změnu statusu. Zůstává nutné pro globálního
Schwarzschildova reprezentanta nebo tvrzení pokrývající všechny prostoročasy
GR. Toto oddělení rozsahů neodstraňuje výše uvedenou lokální překážku výběru
akce.

% BILINGUAL-UNIT: unconditional.verification
\section{Záznam ověření}
Přesný racionální checker
\texttt{tools/verify\_gr\_unconditional\_closure.py} nezávisle ověřuje
rovnici~\eqref{eq:collapse} pro několik nedegenerovaných racionálních tetrádů
a ověřuje, že koeficient křivosti zůstává volným parametrem při všech
zakódovaných kinematických predikátech. Testuje pouze uvedenou algebraickou
implikaci; neodvozuje míru funkcionálního integrálu ani nedokazuje spektrální
stabilitu.

\textbf{LEANEM OVĚŘENÉ POUZE IMPLIKACE Z EXPLICITNÍCH PREMIS.} Lean 4.33.1
s Mathlib 4.33.1 kernelově ověřuje čtyřrozměrnou kontrakci v
\texttt{formal/lean/UBT/GR/CompositeKinetic.lean} a logickou implikaci
neurčenosti koeficientu křivosti v
\texttt{formal/lean/UBT/GR/CurvatureUnderdetermination.lean}. Připnutý Lake
manifest a cíl CI činí tyto omezené kontroly reprodukovatelnými. Druhá věta
zejména předpokládá, místo aby dokazovala, že model s nulovým koeficientem
splňuje kinematický predikát UBT. Žádná z vět nedokazuje výběr ani jednoznačnost
akce a neformalizuje ani neodvozuje mikroskopickou míru funkcionálního
integrálu, kompozitní Hessián nebo fyzikální stabilitu $\psi$. Proto se
nepočítají jako důkaz uzavření žádného otevřeného mikroskopického GR gapu.

% BILINGUAL-UNIT: unconditional.verdict
\section{Verdikt}
\begin{tcolorbox}[colback=yellow!7!white,colframe=yellow!55!black,title=Rozhodnutí o statusu]
\textbf{Status neměnit na vyšší.} Nejsilnějším oprávněným verdiktem zůstává
\textbf{reprodukce GR: CLOSED CONDITIONALLY} pro lokální čtyřrozměrnou
infračervenou Einsteinovu gravitaci. Nepodmíněné \texttt{CLOSED} vyžaduje
odvozený mikroskopický selector, který určí konečný kladný
Einsteinův--Hilbertův koeficient a kontrolovaný kompozitní Hessián/míru, spolu
s fyzikálním výběrem $\psi$-sektoru. Globální slepování leží mimo toto lokální
rozhodnutí.
\end{tcolorbox}

\end{document}
